\begin{scanpage}{0239}{221}
{\cyrillicfont\small
\noindent53. Гринберг Г. А. Новый метод решения задачи дифракции электромагнитных волн на плоскости с безграничной прямолинейной щелью и родственных ей проблем. ЖТФ, 27, № 11, 2595--2605, 1957.\par
54. Гринберг Г. А. Метод решения дифракционных задач для плоских идеально проводящих экранов, основанный на изучении наводимых на экранах теневых токов. ЖТФ, 28, № 3, 542--554, 555--568, 1958.\par
55. Гринберг Г. А. О дифракции электромагнитных волн на полосе конечной ширины. ДАН СССР, 129, № 2, 295--298, 1959.\par
56. Гринберг Г. А. и Пименов Ю. В. К вопросу о дифракции электромагнитных волн на идеально проводящей плоскости с круглым отверстием. Ж. Т. Ф. 29, № 10, 1206--1211, 1959.\par
57. Гринберг Г. А. и Колесникова Э. М. К вопросу о дифракции электромагнитных волн на идеально проводящем плоском кольце. ЖТФ, 31, № 1, 13--17, 1961.\par
58. Millar R. F. Diffraction by a wide slit and complementary strip. Proc. Cambridge Philosophical Society, 54, 1, 479--496, 497--511, 1958.\par
59. Кинг Р. и У Тай-Цзунь. Рассеяние и дифракция электромагнитных волн. Перевод с английского под редакцией Э. Л. Бурштейна. Изд-во иностранной литературы, 1952.\par
60. Кинбер Б. Е. О боковом излучении зеркальных антенн. «Радиотехника и электроника», № 4, 515--558, 1961.\par
61. Кинбер Б. Е. Развязка между близко расположенными зеркальными антеннами. «Радиотехника и электроника», 6, № 6, 907--916, 1961.\par
62. Кинбер Б. Е. О дифракции электромагнитных волн на вогнутой поверхности кругового цилиндра. «Радиотехника и электроника», 6, № 8, 1273--1283, 1961.\par
63. Кинбер Б. Е. О дифракции электромагнитных волн на вогнутой поверхности сферы. «Радиотехника и электроника», 6, № 10, 1652--1657, 1961.\par
64. Фок В. А. Теория дифракции от параболоида вращения. Сборник «Дифракция электромагнитных волн на некоторых телах вращения». Изд-во «Советское радио», 1957.\par
65. Фок В. А. и Федоров А. А. Дифракция плоской электромагнитной волны на идеально проводящем параболоиде вращения. ЖТФ, 28, № 11, 2548--2566, 1958.\par
66. Фок В. А. Поле от вертикального и горизонтального диполя, приподнятого над поверхностью земли. ЖЭТФ, 19, № 10, 916--929, 1949.\par
67. Леонтович М. А. Об одном методе решения задач о распространении электромагнитных волн вдоль поверхности земли. Изв. АН СССР, сер. физ. 8, 16--22, 1944.\par
68. Леонтович М. А. и Фок В. А. Решение задачи о распространении электромагнитных волн вдоль поверхности земли по методу параболического уравнения. ЖЭТФ, 16, № 7, 557--573, 1946.\par
69. Фок В. А. Теория распространения радиоволн в неоднородной атмосфере для приподнятого источника. Изв. АН СССР, 14, № 1, 70--94, 1950.\par
70. Фок В. А. Поле плоской волны вблизи поверхности проводящего тела. Изв. АН СССР, сер. физ. 10, № 2, 171--185, 1946.\par
71. Вайнштейн Л. А., Федоров А. А. Рассеяние плоских и цилиндрических волн на эллиптическом цилиндре и концепция дифракционных лучей. «Радиотехника и электроника», 6, № 1, 31--46, 1961.\par
72. Малюжинец Г. Д. и Вайнштейн Л. А. Поперечная диффузия при дифракции на импедансном цилиндре большого радиуса. ч. I. Параболическое уравнение в лучевых координатах. «Радиотехника и электроника», 6, № 8, 1247--1258, 1961.\par
73. Вайнштейн Л. А., Малюжинец Г. Д. Поперечная диффузия при дифракции на импедансном цилиндре большого радиуса. ч. II. Асимптотические законы дифракции в декартовых координатах. «Радиотехника и электроника», 6, № 9, 1483--1495, 1961.\par
74. Федоров А. А. Асимптотическое решение задачи о дифракции плоской электромагнитной волны на идеально проводящей\par}
\end{scanpage}
